\section{The Architect's Question}\label{the-architects-question}

Having established the single-host serving model --- a 70B FP16 model on
8× H100 serving \textasciitilde2,000 registered users at
\textasciitilde10 rps average / \textasciitilde40 rps peak --- the
natural next question is: when does one host cease to suffice, and how
do we scale to a fleet? This chapter answers the architect's question:
at what point must we move from a lone server to a coordinated fleet,
and what are the quantitative trade-offs of that transition?

\section{Concept}\label{concept}

The transition from a single-server deployment to a fleet is not a
qualitative leap but a quantitative threshold crossing. In the
single-server model, all request routing, load balancing, and health
monitoring occur within one process on one machine. The architect thinks
in terms of request-per-second capacity, GPU memory utilization, and
inference latency on that one host. As traffic grows, the single host
hits limits in three domains: computational throughput (tokens/sec), GPU
memory capacity, and inter-request interference (contention for shared
resources such as NVMe I/O, PCIe bandwidth, and CPU). When any of these
limits is breached, the fleet model introduces multiple identical hosts
behind a load balancer, each independently serving a slice of the total
traffic. The fleet model adds three new architectural ingredients: (a) a
load balancer that distributes requests across hosts, (b) cross-host
communication patterns when a request span multiple machines, and (c)
capacity planning that accounts for redundancy, failover, and the
diminishing returns of horizontal scaling. The key insight is that
horizontal scaling transforms per-host limits into system-wide
capacities, but at the cost of added coordination overhead, state
synchronization, and network budgeting.

\section{Mental Model}\label{mental-model}

We visualize the fleet as a two-level hierarchy. At the top level, a
load balancer --- either hardware (L7 switch) or software (NGINX,
HAProxy, Cloud Load Balancer) --- accepts inbound requests and routes
them to one of N backend hosts. Each backend host runs an identical copy
of the model and serves requests independently. At the bottom level,
each host operates exactly as the single-server model does: a model
engine accepting prompts, producing completions, and returning tokens to
the load balancer, which forwards the response to the client.

The mental model introduces two new invariants. First, the
\textbf{throughput invariant}: total system throughput ≈ N ×
throughput\_per\_host, reduced by a factor ρ that captures load‑balancer
overhead, health‑check traffic, and imperfect distribution (typically ρ
∈ {[}0.85, 0.98{]} depending on LB sophistication). Second, the
\textbf{memory invariant}: total GPU memory across the fleet is N ×
memory\_per\_host, but the per-host memory available for model weights
and KV cache remains the same as the single-host case. The fleet does
not increase the per-host model size; it increases the aggregate number
of concurrent requests the system can serve by parallelizing across
hosts.

\begin{figure}
\centering
\pandocbounded{\includegraphics[keepaspectratio,alt={\hyperref[fig:17.1]{Fig 17.1} - A fleet of one model: a load balancer fans requests out to N identical hosts, each running the same model on its own GPU {[}2° DERIVED{]}},width=\maxwidth{\textwidth},keepaspectratio]{figures/fig-17-1701.pdf}}
\caption{\hyperref[fig:17.1]{Fig 17.1} - A fleet of one model: a load balancer fans requests
out to N identical hosts, each running the same model on its own GPU
{[}2° DERIVED{]}}

\label{fig:17.1}\end{figure}

A critical detail is the \textbf{session-affinity} decision. If user
sessions must maintain conversation state (KV cache, accumulated tool
outputs, retrieval context), the load balancer must route subsequent
requests from the same user to the same host (sticky sessions). This
reduces cross-host communication but limits the effective parallelism,
because a popular user can pin a single host to capacity while other
hosts stand underutilized. If sessions are stateless (the model is
reloaded or KV cache is discarded between requests), any host can serve
any request, and load distribution becomes even. The architect must
choose between density (sticky sessions, higher per-host utilization)
and evenness (stateless routing, lower per-host utilization but better
load balance).

\section{Worked Example}\label{worked-example}

We continue from the canonical scenario: \textasciitilde2,000 registered
users, \textasciitilde5\% concurrent (\textasciitilde100 users),
\textasciitilde10 rps average / \textasciitilde40 rps peak, prompt
\textasciitilde9,200 input tokens + \textasciitilde300 output, 70B-class
dense FP16 model, 8×H100 host with 640\,GB GPU memory. A single 8×H100
host's capacity is set by the \emph{binding constraint} among GPU
compute, HBM bandwidth, and KV-cache memory. Let us verify with
arithmetic, keeping the two quantities distinct: \textbf{aggregate HBM}
(640 GB) is not the same as \textbf{KV-available memory} (640 − 140 GB
weights − runtime/workspace reserve).

\textbf{Single-host concurrency capacity (KV-bound).} The per-token KV
footprint for the canonical 70B model is 2 × layers × hidden × bytes = 2
× 80 × 8,192 × 2 ≈ 2.5 MB/token (\hyperref[chap:7]{Chapter 7} canonical). Subtract the 140
GB FP16 weights and \textasciitilde64 GB runtime/activations/workspace
from the 640 GB pool, leaving ≈436 GB available for KV. Each request at
9,200 input + 300 output reserves 9,500 × 2.5 MB ≈ 23.8 GB of KV. A
single host therefore holds C = 436 GB ÷ 23.8 GB ≈ \textbf{18 concurrent
requests} --- three orders of magnitude below the \textasciitilde440 the
earlier 128-byte/token variant implied. (That variant conflated
aggregate HBM with KV headroom and used a κ \textasciitilde20,000× too
small.)

\textbf{Viability at peak.} Peak traffic is 40 rps; with a target
latency of \textasciitilde1 s that is 40 concurrent requests in flight,
against a KV capacity of \textasciitilde18. \textbf{A single host is not
viable at peak even for the non-agentic canonical workload} --- it would
need \textasciitilde3 hosts at 40 rps, and \textasciitilde2 more for
growth to 100 rps. This flips the original conclusion: the fleet is not
a distant luxury but a necessity at the canonical peak, and it is driven
primarily by KV residency.

\textbf{KV compression is the first lever.} Switching KV to FP8 (≈1.3
MB/token, \hyperref[chap:7]{Chapter 7}) cuts per-request KV to ≈12.4 GB and raises C to ≈35
concurrent per host, roughly doubling headroom before buying hardware.
This is why the architect treats KV precision as a first-class capacity
dial, not a cosmetic detail.

\textbf{Threshold for fleet expansion.} Suppose the agentic layer from
Chapter\,16 is added, with a median of 2 turns and a 95th‑percentile of
4 turns. Each turn adds \textasciitilde800 tokens of retrieved context
and \textasciitilde65 tokens of model-generated reasoning, so the
effective per‑request token count is

\[
I(T) = I_0 + T(\delta + \gamma) + O_\text{final} = 9{,}200 + T \times 865 + 300
\]

\begin{itemize}
\tightlist
\item
  Median (T=2):
  \(I(2) = 9{,}200 + 2 \times 800 + 2 \times 65 + 300 = 10{,}630\)
  tokens
\item
  95th percentile (T=4):
  \(I(4) = 9{,}200 + 4 \times 800 + 4 \times 65 + 300 = 12{,}740\)
  tokens
\end{itemize}

The KV cache per request rises proportionally, and per-host concurrency
capacity follows from the Ch17 formula \(C = M_{kv}/(I(T)\kappa)\):

\begin{itemize}
\tightlist
\item
  at T=2:
  \(10{,}630 \times 2.5 \text{ MB} \approx 26.6 \text{ GB/request} \Rightarrow C \approx 16\)
  concurrent
\item
  at T=4:
  \(12{,}740 \times 2.5 \text{ MB} \approx 31.9 \text{ GB/request} \Rightarrow C \approx 14\)
  concurrent
\end{itemize}

Even the \emph{median} agentic request already drops a single host below
the 40-concurrent peak demand, so the fleet is required under almost any
agentic assumption.

\textbf{Scaling forecast.} To hold the canonical 40 rps at
\textasciitilde1 s latency with FP16 KV (baseline C≈18), the fleet needs
⌈40/18⌉ ≈ 3 hosts (≈4 at the 70\% utilization target the table applies);
at the agentic T=4 (C≈14), peak demand of 40 concurrent needs ⌈40/14⌉ ≈
3 hosts with thin margin (≈5 at the 70\% margin), and growth to 100 rps
(100 concurrent) needs ⌈100/14⌉ ≈ 8 hosts (FP16 KV) --- or
\textasciitilde4--5 hosts if KV is FP8 (C≈26 at T=4). The load balancer
distributes λ/N rps per host and λ·L/N concurrent per host, and each
host's per-host KV usage must stay comfortably under its C. This is the
fleet point: hosts are added because \emph{concurrency demand ×
KV-per-request} exceeds a single host's residency, at the cost of
multiplied capital expenditure and load‑balancer overhead.

\textbf{Cross-host communication.} In the canonical scenario, requests
are independent --- no user session spans multiple hosts. If an agentic
loop required retrieving results from a prior host (e.g., distributed
retrieval), the system would need a coordination layer (Redis, gRPC, or
message queue) with added latency. We quantify this below in the
measurement section.

\section{Measurement}\label{measurement}

We derive three real quantities from the worked example and from
production traffic patterns.

\subsection{Per-host throughput and KV cache
budget}\label{per-host-throughput-and-kv-cache-budget}

Let H be the number of hosts, λ the incoming rps peak, T the random
variable for agent turn count with distribution P(T), δ = 800 tokens per
turn of retrieved context, γ = 65 tokens of model-generated reasoning
per turn, and O\_final = 300 final output tokens. The per-request token
count is

\[
I(T) = I_0 + T \cdot \delta + T \cdot \gamma + O_\text{final}
\]

where I₀ = 9,200 is the baseline input tokens. The per-host
concurrent‑request capacity in terms of KV cache is

\[
C = \frac{M_{kv}}{I(T) \cdot \kappa}
\]

where M\_kv is the KV‑available memory --- the portion of on‑host HBM
left for KV after weights and runtime --- and
\(\kappa = 2 \times n_\text{layers} \times d_\text{hidden} \times \text{bytes} \approx 2.5\)
MB/token is the per‑token KV footprint (\hyperref[chap:7]{Chapter 7} canonical; NOT the
\textasciitilde128 B of an earlier draft, which omitted the layer
count). For the canonical host, M\_kv = 640 GB − 140 GB (FP16 weights) −
\textasciitilde64 GB (runtime/activations/workspace/NCCL) ≈ 436 GB.
Dividing memory by bytes (not the byte‑divided‑by‑bytes collapse that
trips unit analysis):

\begin{itemize}
\tightlist
\item
  For T = 0 (baseline, I = 9,500): C = 436 GB / (9,500 × 2.5 MB) ≈
  436/23.8 ≈ \textbf{18 concurrent requests per host}.
\item
  For T = 2 (median): C = 436 GB / (10,630 × 2.5 MB) ≈ 436/26.6 ≈
  \textbf{16 concurrent requests per host}.
\item
  For T = 4 (95th percentile): C = 436 GB / (12,740 × 2.5 MB) ≈ 436/31.9
  ≈ \textbf{14 concurrent requests per host}.
\end{itemize}

These numbers are the KV-residency ceiling; measured concurrency may be
lower if GPU throughput or bandwidth binds first, so the architect takes
the min of the KV, compute, and bandwidth capacities.

If the peak traffic λ\_peak = 40 rps at an average latency L, the
concurrency in flight is ≈ λ\_peak × L, and the minimum number of hosts
is

\[
H_\text{min} = \left\lceil \frac{\lambda_\text{peak} \cdot L}{C} \cdot \frac{1}{u} \right\rceil
\]

where \(u\) is the utilization target.

\begin{figure}
\centering
\pandocbounded{\includegraphics[keepaspectratio,alt={\hyperref[fig:17.2]{Fig 17.2} --- Host-count contours H = λ·L/(C·util) over peak rate λ × average latency L, with the FP16-KV family (C ≈ 18, filled) and the FP8-KV family (C ≈ 35, dashed) {[}DERIVED: Ch17 §3 + Ch7 KV constants{]}},width=\maxwidth{\textwidth},keepaspectratio]{figures/fig-17-1702.pdf}}
\caption{\hyperref[fig:17.2]{Fig 17.2} --- Host-count contours H = λ·L/(C·util) over peak
rate λ × average latency L, with the FP16-KV family (C ≈ 18, filled) and
the FP8-KV family (C ≈ 35, dashed) {[}DERIVED: Ch17 §3 + Ch7 KV
constants{]}}

\label{fig:17.2}\end{figure}

\emph{\hyperref[fig:17.2]{Fig 17.2} --- The fleet-sizing formula H = ⌈λ·L/(C·util)⌉ rendered
as one glanceable surface. Each contour is a fixed host count; moving to
the upper-right (higher peak rate or higher latency) costs hosts. The
dashed FP8-KV family sits to the right of the FP16 family at every
contour --- a reminder that halving the KV byte width is the cheapest
way to shrink the fleet before buying hardware. The black dots mark the
canonical operating points (Section 3).}

\begin{itemize}
\tightlist
\item
  At T=4 with L=1\,s: concurrency = 40, C ≈ 14 → H\_min = ⌈40/14⌉ ≈ 3
  hosts (≈5 at the 70\% utilization ceiling).
\item
  At the same T=4 with L=2\,s: concurrency = 80, H\_min = ⌈80/14⌉ ≈ 6
  hosts.
\end{itemize}

If traffic grows to λ\_peak = 100 rps:

\begin{itemize}
\tightlist
\item
  At T=4, L=1\,s: concurrency ≈ 100, C ≈ 14 → H\_min = ⌈100/14⌉ ≈ 8
  hosts (FP16 KV).
\item
  Switching to FP8 KV (κ ≈ 1.3 MB/token, C ≈ 26 at T=4) roughly halves
  this to \textasciitilde4--5 hosts --- the KV-precision dial again
  proving more cost-effective than buying hardware.
\end{itemize}

This formulation keeps the two load drivers separate:
\textbf{latency-scaled concurrency} (λ·L) on the demand side and
\textbf{KV residency capacity} (C) on the supply side. Raw rps alone
under-sizes fleets when latency climbs.

\subsection{Load‑balancer overhead
factor}\label{loadbalancer-overhead-factor}

Let ρ be the effective throughput fraction after load‑balancer effects.
Empirically, a well‑configured software LB (NGINX) achieves ρ ≈ 0.95,
meaning total system throughput = ρ · N · throughput\_per\_host. A
hardware L7 switch may achieve ρ ≈ 0.98. The overhead comes from:
health‑check traffic (\textasciitilde1\% of requests), session‑affinity
hashing skew, and LB process CPU limits. The throughput\_per\_host
itself must come from a benchmark on the actual workload
(tokens/context/batch) --- it is an ILLUSTRATIVE proxy until measured.
What is \emph{derivable} is the KV‑residency ceiling C from Section 3;
the LB overhead multiplies the schedulable fraction of that ceiling. In
the worked example, the fleet is currently sized by KV capacity (Section
3), not by a fabricated tokens/s number: at T=4, FP16-KV C≈14/host sets
the ceiling, and adding a host raises aggregate schedulable concurrency
by ρ·C, i.e.~\textasciitilde13 concurrent at ρ=0.95.

\subsection{Cross-host interconnect
budget}\label{cross-host-interconnect-budget}

If a fraction φ of requests require cross-host data (e.g., distributed
retrieval, shared vector index), each such request incurs an additional
network round‑trip latency λ\_net and moves data volume D over the
host‑to‑host network. Assuming an intra‑rack 10\,Gbps Ethernet backbone,
the per-request bandwidth drawn on the network is

\[
B_\text{per-req} = \frac{D}{\lambda_\text{net} \cdot \varphi \cdot N}
\]

For D = 5\,MB (top‑5 retrieval snippets), λ\_net = 20\,ms (typical rack
latency), φ = 0.1 (10\% of requests cross‑host), N = 4:
\(B_\text{per-req} = 5 \text{ MB} / (0.02 \text{ s} \times 0.1 \times 4) \approx 6.25 \text{ MB/s} \approx 50\)\,Mbps,
well within 10\,Gbps capacity. However, if φ scales to 0.5 (many
distributed queries) and D grows to 50\,MB, the per-request bandwidth
becomes \(50/(0.02 \times 0.5 \times 4) = 12.5\)\,MB/s ≈ 100\,Mbps,
still fine, but aggregate traffic across the backbone is

\[
\text{aggregate} = \varphi \cdot N \cdot \frac{D}{\lambda_\text{net}} = 0.5 \times 4 \times \frac{50 \text{ MB}}{0.02 \text{ s}} = 5 \text{ GB/s}
\]

which exceeds the 10 Gbps (\textasciitilde1.25\,GB/s) limit and would
require 5× overprovisioning or a wider fabric (e.g., 50\,Gbps or
100\,Gbps). This calculation shows that cross-host communication must be
dimensioned early, even if the base case appears benign.

\section{Common Mistakes}\label{common-mistakes}

\begin{enumerate}
\def\labelenumi{\arabic{enumi}.}
\item
  \textbf{Ignoring KV cache growth from agentic loops.} The most
  frequent architectual error is to assume that adding an agentic layer
  barely changes per-request resource usage. As shown, each turn adds
  \textasciitilde800 tokens of context, which linearly reduces
  concurrent‑request capacity. Forgetting this leads to
  under‑provisioning the fleet.
\item
  \textbf{Over‑relying on sticky sessions without load‑balance
  analysis.} Sticky sessions simplify cross‑host communication but can
  create hot‑spot hosts if a few users generate disproportionate
  traffic. Always measure the session‑affinity distribution; if the Gini
  coefficient of per-host request counts exceeds 0.3, consider stateless
  routing or session sharding.
\item
  \textbf{Using raw rps to size the fleet.} Requests‑per‑second is
  necessary but not sufficient; the product of rps × average latency
  gives concurrent load, which must fit within KV cache capacity. A host
  serving 40 rps at 0.5\,s latency carries 20 concurrent requests,
  whereas at 2\,s latency it carries 80 --- the same rps, very different
  resource usage.
\item
  \textbf{Under‑estimating load‑balancer overhead.} A 5\% throughput
  reduction from ρ = 0.95 may seem small, but at the margin of capacity
  it can be the difference between fitting on N hosts and needing N+1.
  Always provision with ρ explicitly in the formula.
\item
  \textbf{Failing to model latency tail under spike traffic.} GPU
  inference time is not constant; under sustained high load, kernel‑mode
  preemption, memory fragmentation, and concurrent kernel launches can
  increase per‑token latency by 20--50\%. The linear model L =
  λ\_inference + λ\_tool is a lower bound; the architect should add a
  20\% safety margin.
\end{enumerate}

\section{Architecture Consequence}\label{architecture-consequence}

Transitioning from one host to a fleet reshapes every layer of the
system.

\textbf{GPU fleet sizing.} The number of hosts N must satisfy the
latency‑augmented capacity formula H = ⌈ λ\_peak · L / (C ·
util\_target) ⌉, where C itself depends on the agentic turn
distribution. If the team expects T ≤ 3 for 90\% of requests and λ\_peak
= 100 rps with L ≈ 2\,s (inference + agentic overhead), the concurrency
in flight is 200. With FP16 KV at T=3, C ≈ 14, so H = ⌈200/14⌉ ≈ 15
hosts --- a large fleet. Switching to FP8 KV (κ ≈ 1.3 MB/token → C ≈ 28
at T=3) roughly halves this to \textasciitilde8 hosts. Either way, the
KV‑residency capacity --- not raw rps --- is the binding term, and the
honest answer is a multi‑host fleet, not a near‑single‑host illusion.
The key takeaway: size the fleet for the 95th‑percentile latency and
token scenario, not the median, and treat KV precision as the cheapest
way to shrink the fleet before adding hardware.

\textbf{Load‑balancer selection and configuration.} The choice between
hardware (F5, Cloud LB) and software (NGINX, HAProxy) affects ρ and the
session‑affinity model. Hardware L7 switches offer better ρ (≈0.98) but
less fine‑grained traffic‑shaping; software LBs offer ρ ≈ 0.95 but can
implement consistent hashing for session affinity. The architect should
match the LB capability to the session model: sticky sessions → any LB;
stateless → software LB with consistent hashing.

\textbf{Cross-host data infrastructure.} If any request type requires
cross‑host communication (distributed retrieval, multi‑model
orchestration, state sync), a dedicated backend network must be
provisioned. The minimum bandwidth is B\_min = φ · N · D / λ\_net, and
the fabric should provide at least 2× B\_min for headroom. In the
canonical scenario with φ = 0.1, N = 4, D = 5\,MB, λ\_net = 20\,ms,
B\_min ≈ 0.5\,Gbps, so a 10\,Gbps intra‑rack network is ample. But the
architect should instrument φ and D from production telemetry and
revisit the bandwidth provisioning quarterly.

\textbf{Observability across the fleet.} With N hosts, tracing a request
end‑to‑end requires correlating the LB routing decision, the selected
host's inference logs, and any cross‑host data calls. Structured logging
with a request‑ID field that propagates from the client through the LB
to the host and back is essential. Distributed tracing (OpenTelemetry,
Jaeger) should be enabled by default. Alerts should trigger when
per‑host concurrency exceeds 80\% of C, when LB error rate exceeds
0.1\%, or when cross‑host bandwidth utilization exceeds 70\%.

\textbf{Failover and redundancy.} With N ≥ 2, the fleet provides
automatic failover: if one host crashes, the LB routes its traffic to
the remaining N−1 hosts, reducing capacity by \textasciitilde1/(N−1).
For N = 4, a single‑host loss reduces capacity from 100\% to 75\%
(assuming equal load distribution). The architect must decide whether
this is acceptable or whether N must be larger (e.g., N = 6 for 83\%
residual capacity after one failure). Health‑check frequency and
graceful drain procedures affect the real-world failover time, typically
5--15 seconds.

\section{What We Still Don't Know}\label{what-we-still-dont-know}

Despite the arithmetic above, several questions remain open and would
benefit from targeted research:

\begin{itemize}
\tightlist
\item
  \textbf{Turn-count distribution under fleet load.} Do larger fleets (N
  \textgreater{} 4) change the observed T distribution? Intuitively,
  more hosts reduce per-host contention, potentially lowering latency
  and reducing the agentic loop depth, but empirical measurement is
  needed.
\item
  \textbf{Latency amplification non-linearity.} Our linear model β(T) =
  T·L / L₀ assumes constant per‑turn latency. Under fleet scaling, GPU
  saturation effects may make L increase with concurrent requests,
  introducing convexity. Empirical load‑testing across N = 1, 2, 4, 8
  hosts is required.
\item
  \textbf{Optimal fleet size for cost vs.~latency.} The cost of adding a
  host is roughly the cost of 8×H100 (multi‑hundred-thousand dollars).
  The marginal latency improvement from the Nth host diminishes
  (diminishing returns). A cost‑per‑latency‑improved curve would help
  the architect justify fleet size to stakeholders.
\item
  \textbf{Cross‑host communication patterns at scale.} The φ = 0.1
  assumption may not hold for all workloads (e.g., RAG over a shared
  index, distributed fine‑tuning). Real‑world measurement of cross‑host
  traffic patterns is needed to size the backbone accurately.
\item
  \textbf{Context‑caching efficacy across a fleet.} If many users query
  related topics, can a fleet‑wide KV cache or retrieval cache reduce
  the effective δ per request? The savings could be substantial but
  require cross-host cache-invalidation protocols that are not yet well
  understood.
\end{itemize}

\section{End-of-Chapter Mini-Case}\label{end-of-chapter-mini-case}

\textbf{Scenario.} A AI‑product company starts with a single 8×H100 host
serving the canonical traffic of 2,000 users at 10 rps average / 40 rps
peak, 9,200 input + 300 output tokens, no agentic layer. After six
months, user growth pushes peak traffic to 100 rps, and the team
introduces the agentic layer from Chapter\,16 with a maximum of 3 turns
(median T = 2, 95th‑percentile T = 4). The team must decide on fleet
size.

\textbf{Initial sizing analysis.} At T = 3 (worst‑case assumed for
sizing), I(3) = 9,200 + 3·800 + 3·65 + 300 = 12,095 tokens. With FP16 KV
(κ = 2.5 MB/token), the per‑request KV is 12,095 × 2.5 MB ≈ 30.2 GB, so
KV‑residency capacity per host is C = M\_kv ÷ 30.2 GB. With M\_kv ≈ 436
GB (640 − 140 weights − \textasciitilde64 runtime), C ≈ 14 concurrent
requests per host. With peak latency L = 2\,s (inference + 3 agent
turns), peak concurrency = λ\_peak × L = 100 × 2 = 200. The fleet must
therefore satisfy 200 concurrent in flight against \textasciitilde14 per
host: even at 70\% utilization target, H = ⌈200/(14 × 0.7)⌉ ≈ 21 hosts
on FP16 KV --- a much larger fleet than a naive rps-based estimate would
suggest. This is exactly why the KV‑residency model matters: it turns a
plausible-looking ``1--5 hosts'' into an honest ``10--20+ hosts''.

\textbf{KV compression changes the answer.} Switching the KV cache to
FP8 (κ ≈ 1.3 MB/token) cuts per‑request KV at T=3 to 12,095 × 1.3 MB ≈
15.7 GB and roughly doubles C to ≈28 per host. The fleet then needs H =
⌈200/(28 × 0.7)⌉ ≈ 11 hosts. Combined with prompt/prefix caching (which
can eliminate the repeated 9,200‑token base prompt from the KV per
turn), the practical fleet lands in the 6--12 host range rather than 21.
KV precision and prefix caching are the primary levers that make agentic
fleets affordable.

\textbf{Decision.} The team procures \textasciitilde12 hosts (an 8--16
host range covering the FP8+prefix-cache sizing), runs NGINX as a
software load balancer with consistent hashing for session affinity, and
instruments each host with OpenTelemetry tracing. Cross‑host
communication is minimal (φ ≈ 0.05 because most queries are standalone),
so the existing 10\,Gbps intra‑rack network suffices. Monthly
GPU‑utilization metrics should run in the 55--70\% band against the
KV‑residency ceiling, and the SLA of \textless1\,s 95th‑percentile
latency must be validated by a deployment benchmark (\hyperref[chap:14]{Chapter 14}) ---
end‑to‑end latency including agentic overhead is an ILLUSTRATIVE target
until measured.

\textbf{Lesson.} The fleet decision is driven by KV‑residency capacity
vs.~latency‑scaled concurrency --- not by raw rps, and not by a
per‑layer KV constant. The architect's checklist: (1) compute I(T) for
the expected turn distribution, (2) derive per‑host concurrent capacity
C from the \emph{canonical} per‑token KV (2 × layers × hidden × bytes),
(3) convert peak rps into concurrency via λ·L, (4) compute H =
⌈λ·L/(C·util\_target)⌉, (5) apply the cheapest capacity levers first ---
KV precision (FP8), prefix caching, attention architecture (GQA/MQA) ---
before buying hosts, and (6) validate throughput and latency with a
benchmark, not an assumption. A 12‑host FP8 fleet costs meaningful GPU
capital, but avoids both wasteful over‑provisioning and a costly
re‑platform when traffic doubles.

\begin{center}\rule{0.5\linewidth}{0.5pt}\end{center}

\textbf{\hyperref[tab:17.1]{Table 17-1}: Fleet sizing matrix --- KV-residency ceiling and
minimum hosts (canonical FP16 KV).}

C (per-host concurrent capacity, from canonical 2.5 MB/token and
M\_kv≈436 GB): baseline ≈18, T=2 ≈16, T=4 ≈14. Concurrency in flight = λ
× L. Minimum hosts H = ⌈λ·L / (C·0.7)⌉ at a 70\% utilization ceiling.

\begin{longtable}[]{@{}
  >{\raggedright\arraybackslash}p{0.2000\linewidth}
  >{\raggedright\arraybackslash}p{0.1926\linewidth}
  >{\raggedright\arraybackslash}p{0.1926\linewidth}
  >{\raggedright\arraybackslash}p{0.1185\linewidth}
  >{\raggedright\arraybackslash}p{0.1185\linewidth}
  >{\raggedright\arraybackslash}p{0.1778\linewidth}@{}}
\toprule\noalign{}
\begin{minipage}[b]{\linewidth}\raggedright
Traffic scenario (λ peak)
\end{minipage} & \begin{minipage}[b]{\linewidth}\raggedright
L=1s, T=0 → concurrency
\end{minipage} & \begin{minipage}[b]{\linewidth}\raggedright
L=1s, T=4 → concurrency
\end{minipage} & \begin{minipage}[b]{\linewidth}\raggedright
H @ T=0 (C=18)
\end{minipage} & \begin{minipage}[b]{\linewidth}\raggedright
H @ T=4 (C=14)
\end{minipage} & \begin{minipage}[b]{\linewidth}\raggedright
H @ T=4, FP8 KV (C≈26)
\end{minipage} \\
\midrule\noalign{}
\endhead
\bottomrule\noalign{}
\endlastfoot
40 rps (canonical peak) & 40 & 40 & ⌈40/12.6⌉ ≈ 4 & ⌈40/9.8⌉ ≈ 5 & ≈
3 \\
100 rps (growth) & 100 & 100 & ⌈100/12.6⌉ ≈ 8 & ⌈100/9.8⌉ ≈ 11 & ≈ 6 \\
200 rps (2× growth) & 200 & 200 & ⌈200/12.6⌉ ≈ 16 & ⌈200/9.8⌉ ≈ 21 & ≈
11 \\

\label{tab:17.1}\end{longtable}

\emph{C is derived (canonical KV); λ, L, and the 70\% utilization target
are ILLUSTRATIVE scenario inputs to be validated by a deployment
benchmark (\hyperref[chap:14]{Chapter 14}) before committing hardware. H is a KV-residency
lower bound --- the true answer is the max of this KV bound and the
measured throughput bound.}

\begin{center}\rule{0.5\linewidth}{0.5pt}\end{center}
